% source stamp shared:2ed4dbc6 — maintained by build-source-stamps.py; it is what makes editing ../macros.sty, the pictures or the book's numbering recompile every block
% Each extraction job runs in html/.reflowtex-build/HASH, i.e. always exactly
% three directories below the book repository root, so ../../../ reaches the
% book's own LaTeX sources (macros.sty, knowledges.tex, main.aux) wherever the
% repository happens to be checked out. There is no second copy of the sources
% for the web edition: this is the same macros.sty the PDF build uses. Never
% write an absolute, machine-specific path here.
\usepackage{../../../macros}
\input{../../../knowledges}
\newcommand{\map}{\mathsf{map}}
\usepackage{xr-hyper}
% …against a snapshot of main.aux rather than the file itself: the author's
% editor rewrites that file whenever it rebuilds the PDF, and a block that reads
% it mid-write brings the build down. rebuild.sh makes the copy; ../../ is the
% site root from a block's build directory.
\externaldocument{../../.aux-snapshot/main}

% Keep Inconsolata for \ttfamily. libertine loads fontspec (so the document is
% TU-encoded), but the inconsolata package only ever defines the 8-bit family
% zi4, which has no TU shape — NFSS then silently substitutes the roman face and
% every listing comes out in a proportional serif. Selecting the same typeface
% through fontspec, from the OTFs TeX Live ships alongside the Type1 ones, keeps
% the author's choice and gives \ttfamily somewhere to land.
\setmonofont{Inconsolatazi4}[
  Extension = .otf,
  UprightFont = *-Regular,
  BoldFont = *-Bold,
  Scale = MatchLowercase]

% ── the bibliography ────────────────────────────────────────────────────────
% Each chapter is extracted on its own, and the pipeline runs lualatex only —
% never biber (src/encode/pipeline.py). So no \jobname.bbl can ever exist here,
% however many passes it makes, and every \cite comes out undefined.
%
% In a footnote that is worse than it sounds. macros.sty makes \cite mean
% \incite inside a footnote, and \incite prints author, title and year straight
% from the bibliography data — so with no data it prints *nothing*, and the
% sentence simply stops ("first appear in .").
%
% The book's own bibliography is already resolved by biber during the PDF build,
% next door in ../../../main.bbl, and main.tex \inputs every chapter, so that
% file holds every entry any chapter can cite. This is biblatex's own loader
% with two changes: it reads that file instead of \jobname.bbl, and it drops the
% \blx@ifsigned guard, which would otherwise refuse a .bbl belonging to another
% job. Everything after the input — the toggle, the checksum, the rerun check —
% is biblatex's, left alone, so its bookkeeping stays consistent.
\makeatletter
\def\blx@bblfile{%
  \blx@secinit
  \begingroup
  \blx@bblstart
  \InputIfFileExists{../../../main.bbl}
    {\global\toggletrue{blx@bbldone}%
     \blx@generate@bbl@mdfivesum@found{../../../main.bbl}}
    {\blx@warning@noline{No ../../../main.bbl -- run latexmk on the book first,
       or citations will be empty}%
     \blx@generate@bbl@mdfivesum@notfound}%
  \blx@bblend
  \endgroup
  \blx@check@bbl@rerun@mdfivesum
  \csnumgdef{blx@labelnumber@\the\c@refsection}{0}}
\makeatother

% Citations, clickable. A \cite leaves no link of its own that survives into the
% web edition: hyperref points it at the bibliography, which an isolated chapter
% does not contain, and reflowtex drops a reference it cannot resolve. What it
% does record is a URL written out in full, so each citation label is wrapped in
% one — through biblatex's own `bibhyperref` format, which is applied once per
% key, so a multi-key citation gets one link per key rather than one for the lot.
%
% The scheme is our own: nothing navigates to `cite:…`, and the page intercepts
% the click and shows the entry in the third column (layouts/partials/lean-pane
% .html, showReference). A fragment like `#cite-key` cannot be used instead —
% hyperref splits an href at the `#` and routes it down its internal-link path,
% where it is lost.
\DeclareFieldFormat{bibhyperref}{\href{cite:\thefield{entrykey}}{#1}}

% Inside a footnote, macros.sty makes \cite mean \incite, which prints the names,
% the title and the year itself rather than a label — and prints them with
% \printnames and \printfield, which never pass through `bibhyperref` above. Two
% citations in three are in footnotes, so they would all have stayed dead. This
% is macros.sty's own definition with the printed run wrapped in the same link;
% biblatex's punctuation buffer survives the extra group, so the footnote still
% reads "…Elgot and J. E. Mezei, “On relations defined…”, 1965, and also…".
%
% Kept here rather than in macros.sty because it is the web edition's business:
% the PDF has a bibliography of its own to jump to. If \incite is ever changed
% over there, change it here too.
\DeclareCiteCommand{\incite}
  {\usebibmacro{prenote}}
  {\href{cite:\thefield{entrykey}}{%
     \printnames[given-family]{labelname}%
     \setunit{\addcomma\space}%
     \printfield[citetitle]{title}%
     \setunit{\addcomma\space}%
     \printfield{year}}}
  {\multicitedelim}
  {\usebibmacro{postnote}}

